\section{晚清改革可跳转事件节点}
以下锚点连接灾荒、边疆、甲午、教案、新政、学校、法律和宪政。每一制度节点都保留其核心文书与实施范围的区别。
\subsection{1870—1879}
\eventanchor{EQ-1875-GUANGXU-ACCESSION}{光绪帝即位，慈禧、慈安继续掌握关…}光绪元年（1875年），光绪帝即位，慈禧、慈安继续掌握关键政治影响。核心取证：《清德宗实录》光绪元年相关条；宫中朱批奏折、内阁大库题本与《清史稿》相关志传。此类文件确认机构作出或收到的行政动作，不能跳过地方执行与反馈。来源保留：以光绪元年（1875年）的同期文书为定位起点；官修记录、奏报或外交文本服务特定机构立场，具体日期、规模、地方执行与对手视角须以独立材料复核。
\eventanchor{EQ-1875-YUNNAN-MUSLIM-WAR}{杜文秀政权覆亡前后，云南回民战争…}光绪元年（1875年），杜文秀政权覆亡前后，云南回民战争进入惨烈收束阶段。核心取证：《清德宗实录》光绪元年相关条；军机处档、战事方略及地方奏报。编年证词定位国家叙述中的时间，却需同地方或对方材料校验范围。来源保留：以光绪元年（1875年）的同期文书为定位起点；官修记录、奏报或外交文本服务特定机构立场，具体日期、规模、地方执行与对手视角须以独立材料复核。
\eventanchor{EQ-1876-CHEFOO-CONVENTION}{《烟台条约》签订，扩大英国在华通…}光绪二年（1876年），《烟台条约》签订，扩大英国在华通商和内地活动权利。核心取证：《清德宗实录》光绪二年相关条；《筹办夷务始末》相关年卷与中外照会、条约文本。此类文件确认机构作出或收到的行政动作，不能跳过地方执行与反馈。来源保留：以光绪二年（1876年）的同期文书为定位起点；官修记录、奏报或外交文本服务特定机构立场，具体日期、规模、地方执行与对手视角须以独立材料复核。
\eventanchor{EQ-1876-NORTH-CHINA-FAMINE}{丁戊奇荒在华北及西北造成广泛饥荒…}光绪二年（1876年），丁戊奇荒在华北及西北造成广泛饥荒、迁徙和救济行动。核心取证：《清德宗实录》光绪二年相关条；地方志、督抚奏折及地方档案。编年证词定位国家叙述中的时间，却需同地方或对方材料校验范围。来源保留：以光绪二年（1876年）的同期文书为定位起点；官修记录、奏报或外交文本服务特定机构立场，具体日期、规模、地方执行与对手视角须以独立材料复核。
\eventanchor{EQ-1877-FAMINE-RELIEF-NETWORKS}{赈济网络在灾区展开，官办、绅办和…}光绪三年（1877年），赈济网络在灾区展开，官办、绅办和传教团体的救助并存。核心取证：《清德宗实录》光绪三年相关条；地方志、督抚奏折及地方档案。编年证词定位国家叙述中的时间，却需同地方或对方材料校验范围。来源保留：以光绪三年（1877年）的同期文书为定位起点；官修记录、奏报或外交文本服务特定机构立场，具体日期、规模、地方执行与对手视角须以独立材料复核。
\eventanchor{EQ-1878-Zuo-ZONGTANG-XINJIANG}{左宗棠军收复新疆大部，阿古柏政权…}光绪四年（1878年），左宗棠军收复新疆大部，阿古柏政权崩解。核心取证：《清德宗实录》光绪四年相关条；军机处档、战事方略及地方奏报。编年证词定位国家叙述中的时间，却需同地方或对方材料校验范围。来源保留：以光绪四年（1878年）的同期文书为定位起点；官修记录、奏报或外交文本服务特定机构立场，具体日期、规模、地方执行与对手视角须以独立材料复核。
\eventanchor{EQ-1878-XINJIANG-PROVINCE-DEBATE}{朝廷讨论新疆设省、军政与财政安排}光绪四年（1878年），朝廷讨论新疆设省、军政与财政安排。核心取证：《清德宗实录》光绪四年相关条；宫中朱批奏折、内阁大库题本与《清史稿》相关志传。此类文件确认机构作出或收到的行政动作，不能跳过地方执行与反馈。来源保留：以光绪四年（1878年）的同期文书为定位起点；官修记录、奏报或外交文本服务特定机构立场，具体日期、规模、地方执行与对手视角须以独立材料复核。
\eventanchor{EQ-1879-RYUKYU-ANNEXATION}{日本吞并琉球，清廷的交涉未能恢复…}光绪五年（1879年），日本吞并琉球，清廷的交涉未能恢复旧有册封关系。核心取证：《清德宗实录》光绪五年相关条；《筹办夷务始末》相关年卷与中外照会、条约文本。此类文件确认机构作出或收到的行政动作，不能跳过地方执行与反馈。来源保留：以光绪五年（1879年）的同期文书为定位起点；官修记录、奏报或外交文本服务特定机构立场，具体日期、规模、地方执行与对手视角须以独立材料复核。
\subsection{1880—1889}
\eventanchor{EQ-1880-ILI-NEGOTIATIONS}{清俄围绕伊犁归还和边界条件持续谈…}光绪六年（1880年），清俄围绕伊犁归还和边界条件持续谈判。核心取证：《清德宗实录》光绪六年相关条；《筹办夷务始末》相关年卷与中外照会、条约文本。此类文件确认机构作出或收到的行政动作，不能跳过地方执行与反馈。来源保留：以光绪六年（1880年）的同期文书为定位起点；官修记录、奏报或外交文本服务特定机构立场，具体日期、规模、地方执行与对手视角须以独立材料复核。
\eventanchor{EQ-1881-SAINT-PETERSBURG-TREATY}{《中俄伊犁条约》签订，伊犁大部归…}光绪七年（1881年），《中俄伊犁条约》签订，伊犁大部归还清廷并附带领土、赔款和贸易条件。核心取证：《清德宗实录》光绪七年相关条；《筹办夷务始末》相关年卷与中外照会、条约文本。此类文件确认机构作出或收到的行政动作，不能跳过地方执行与反馈。来源保留：以光绪七年（1881年）的同期文书为定位起点；官修记录、奏报或外交文本服务特定机构立场，具体日期、规模、地方执行与对手视角须以独立材料复核。
\eventanchor{EQ-1882-IMO-INCIDENT}{朝鲜壬午兵变后清军入朝，日本也扩…}光绪八年（1882年），朝鲜壬午兵变后清军入朝，日本也扩大影响。核心取证：《清德宗实录》光绪八年相关条；《筹办夷务始末》相关年卷与中外照会、条约文本。此类文件确认机构作出或收到的行政动作，不能跳过地方执行与反馈。来源保留：以光绪八年（1882年）的同期文书为定位起点；官修记录、奏报或外交文本服务特定机构立场，具体日期、规模、地方执行与对手视角须以独立材料复核。
\eventanchor{EQ-1883-SINO-FRENCH-WAR-BEGIN}{中法因越南北部和宗藩关系冲突升级…}光绪九年（1883年），中法因越南北部和宗藩关系冲突升级为战争。核心取证：《清德宗实录》光绪九年相关条；军机处档、战事方略及地方奏报。编年证词定位国家叙述中的时间，却需同地方或对方材料校验范围。来源保留：以光绪九年（1883年）的同期文书为定位起点；官修记录、奏报或外交文本服务特定机构立场，具体日期、规模、地方执行与对手视角须以独立材料复核。
\eventanchor{EQ-1884-FUZHOU-NAVAL-BATTLE}{马江海战中福建水师遭法国舰队重创}光绪十年（1884年），马江海战中福建水师遭法国舰队重创。核心取证：《清德宗实录》光绪十年相关条；军机处档、战事方略及地方奏报。编年证词定位国家叙述中的时间，却需同地方或对方材料校验范围。来源保留：以光绪十年（1884年）的同期文书为定位起点；官修记录、奏报或外交文本服务特定机构立场，具体日期、规模、地方执行与对手视角须以独立材料复核。
\eventanchor{EQ-1884-TAIWAN-BLOCKADE}{法军进攻台湾北部并封锁，台湾成为…}光绪十年（1884年），法军进攻台湾北部并封锁，台湾成为中法战争重要战场。核心取证：《清德宗实录》光绪十年相关条；军机处档、战事方略及地方奏报。编年证词定位国家叙述中的时间，却需同地方或对方材料校验范围。来源保留：以光绪十年（1884年）的同期文书为定位起点；官修记录、奏报或外交文本服务特定机构立场，具体日期、规模、地方执行与对手视角须以独立材料复核。
\eventanchor{EQ-1885-TIANJIN-CHINA-FRANCE-TREATY}{《中法新约》确认越南脱离清朝宗藩…}光绪十一年（1885年），《中法新约》确认越南脱离清朝宗藩体系，战争结束。核心取证：《清德宗实录》光绪十一年相关条；《筹办夷务始末》相关年卷与中外照会、条约文本。此类文件确认机构作出或收到的行政动作，不能跳过地方执行与反馈。来源保留：以光绪十一年（1885年）的同期文书为定位起点；官修记录、奏报或外交文本服务特定机构立场，具体日期、规模、地方执行与对手视角须以独立材料复核。
\eventanchor{EQ-1885-TAIWAN-PROVINCE}{清廷决定将台湾建省，加强行政、军…}光绪十一年（1885年），清廷决定将台湾建省，加强行政、军防和基础设施经营。核心取证：《清德宗实录》光绪十一年相关条；宫中朱批奏折、内阁大库题本与《清史稿》相关志传。此类文件确认机构作出或收到的行政动作，不能跳过地方执行与反馈。来源保留：以光绪十一年（1885年）的同期文书为定位起点；官修记录、奏报或外交文本服务特定机构立场，具体日期、规模、地方执行与对手视角须以独立材料复核。
\eventanchor{EQ-1886-BEIYANG-NAVY-EXPANSION}{北洋海军和沿海防务建设持续推进}光绪十二年（1886年），北洋海军和沿海防务建设持续推进。核心取证：《清德宗实录》光绪十二年相关条；军机处档、战事方略及地方奏报。编年证词定位国家叙述中的时间，却需同地方或对方材料校验范围。来源保留：以光绪十二年（1886年）的同期文书为定位起点；官修记录、奏报或外交文本服务特定机构立场，具体日期、规模、地方执行与对手视角须以独立材料复核。
\eventanchor{EQ-1887-TAIWAN-RAILWAY-PLANNING}{台湾巡抚刘铭传推动铁路、电报、矿…}光绪十三年（1887年），台湾巡抚刘铭传推动铁路、电报、矿务和新式行政项目。核心取证：《清德宗实录》光绪十三年相关条；户部奏销、盐政、河工或海关档案。此类文件确认机构作出或收到的行政动作，不能跳过地方执行与反馈。来源保留：以光绪十三年（1887年）的同期文书为定位起点；官修记录、奏报或外交文本服务特定机构立场，具体日期、规模、地方执行与对手视角须以独立材料复核。
\eventanchor{EQ-1888-NAVAL-AND-ARMY-REFORM}{清廷讨论海军、陆军训练和军工体系…}光绪十四年（1888年），清廷讨论海军、陆军训练和军工体系的经费与制度问题。核心取证：《清德宗实录》光绪十四年相关条；宫中朱批奏折、内阁大库题本与《清史稿》相关志传。此类文件确认机构作出或收到的行政动作，不能跳过地方执行与反馈。来源保留：以光绪十四年（1888年）的同期文书为定位起点；官修记录、奏报或外交文本服务特定机构立场，具体日期、规模、地方执行与对手视角须以独立材料复核。
\eventanchor{EQ-1889-GUANGXU-MARRIAGE}{光绪帝大婚并开始形式上的亲政安排}光绪十五年（1889年），光绪帝大婚并开始形式上的亲政安排。核心取证：《清德宗实录》光绪十五年相关条；宫中朱批奏折、内阁大库题本与《清史稿》相关志传。此类文件确认机构作出或收到的行政动作，不能跳过地方执行与反馈。来源保留：以光绪十五年（1889年）的同期文书为定位起点；官修记录、奏报或外交文本服务特定机构立场，具体日期、规模、地方执行与对手视角须以独立材料复核。
\subsection{1890—1899}
\eventanchor{EQ-1890-KOREAN-POLITICAL-CONFLICT}{朝鲜内政纷争持续牵动清日两国在半…}光绪十六年（1890年），朝鲜内政纷争持续牵动清日两国在半岛的影响力竞争。核心取证：《清德宗实录》光绪十六年相关条；《筹办夷务始末》相关年卷与中外照会、条约文本。此类文件确认机构作出或收到的行政动作，不能跳过地方执行与反馈。来源保留：以光绪十六年（1890年）的同期文书为定位起点；官修记录、奏报或外交文本服务特定机构立场，具体日期、规模、地方执行与对手视角须以独立材料复核。
\eventanchor{EQ-1891-MISSIONARY-CONFLICTS}{各地教案和宗教纠纷继续引发地方治…}光绪十七年（1891年），各地教案和宗教纠纷继续引发地方治理与外交压力。核心取证：《清德宗实录》光绪十七年相关条；地方志、督抚奏折及地方档案。编年证词定位国家叙述中的时间，却需同地方或对方材料校验范围。来源保留：以光绪十七年（1891年）的同期文书为定位起点；官修记录、奏报或外交文本服务特定机构立场，具体日期、规模、地方执行与对手视角须以独立材料复核。
\eventanchor{EQ-1892-RAILWAY-DEBATE}{官员、商人和士绅围绕铁路修筑、主…}光绪十八年（1892年），官员、商人和士绅围绕铁路修筑、主权与资金展开争论。核心取证：《清德宗实录》光绪十八年相关条；户部奏销、盐政、河工或海关档案。此类文件确认机构作出或收到的行政动作，不能跳过地方执行与反馈。来源保留：以光绪十八年（1892年）的同期文书为定位起点；官修记录、奏报或外交文本服务特定机构立场，具体日期、规模、地方执行与对手视角须以独立材料复核。
\eventanchor{EQ-1893-KOREAN-REFORM-CRISIS}{朝鲜改革与政治危机使清日关系持续…}光绪十九年（1893年），朝鲜改革与政治危机使清日关系持续紧张。核心取证：《清德宗实录》光绪十九年相关条；《筹办夷务始末》相关年卷与中外照会、条约文本。此类文件确认机构作出或收到的行政动作，不能跳过地方执行与反馈。来源保留：以光绪十九年（1893年）的同期文书为定位起点；官修记录、奏报或外交文本服务特定机构立场，具体日期、规模、地方执行与对手视角须以独立材料复核。
\eventanchor{EQ-1894-DONGHAK-PEASANT-WAR}{东学农民战争促使清日两国派兵入朝…}光绪二十年（1894年），东学农民战争促使清日两国派兵入朝，危机迅速国际化。核心取证：《清德宗实录》光绪二十年相关条；军机处档、战事方略及地方奏报。编年证词定位国家叙述中的时间，却需同地方或对方材料校验范围。来源保留：以光绪二十年（1894年）的同期文书为定位起点；官修记录、奏报或外交文本服务特定机构立场，具体日期、规模、地方执行与对手视角须以独立材料复核。
\eventanchor{EQ-1894-FIRST-SINO-JAPANESE-WAR}{甲午战争爆发，清日围绕朝鲜和区域…}光绪二十年（1894年），甲午战争爆发，清日围绕朝鲜和区域秩序展开海陆战。核心取证：《清德宗实录》光绪二十年相关条；军机处档、战事方略及地方奏报。编年证词定位国家叙述中的时间，却需同地方或对方材料校验范围。来源保留：以光绪二十年（1894年）的同期文书为定位起点；官修记录、奏报或外交文本服务特定机构立场，具体日期、规模、地方执行与对手视角须以独立材料复核。
\eventanchor{EQ-1894-YALU-NAVAL-BATTLE}{黄海海战后北洋舰队受创，海军力量…}光绪二十年（1894年），黄海海战后北洋舰队受创，海军力量对比进一步恶化。核心取证：《清德宗实录》光绪二十年相关条；军机处档、战事方略及地方奏报。编年证词定位国家叙述中的时间，却需同地方或对方材料校验范围。来源保留：以光绪二十年（1894年）的同期文书为定位起点；官修记录、奏报或外交文本服务特定机构立场，具体日期、规模、地方执行与对手视角须以独立材料复核。
\eventanchor{EQ-1895-SHIMONOSEKI-TREATY}{《马关条约》签订，清割让台湾、澎…}光绪二十一年（1895年），《马关条约》签订，清割让台湾、澎湖并承认朝鲜独立，赔款与通商条件加重。核心取证：《清德宗实录》光绪二十一年相关条；《筹办夷务始末》相关年卷与中外照会、条约文本。此类文件确认机构作出或收到的行政动作，不能跳过地方执行与反馈。来源保留：以光绪二十一年（1895年）的同期文书为定位起点；官修记录、奏报或外交文本服务特定机构立场，具体日期、规模、地方执行与对手视角须以独立材料复核。
\eventanchor{EQ-1895-TAIWAN-REPUBLIC-RESISTANCE}{台湾割让后地方力量建立台湾民主国…}光绪二十一年（1895年），台湾割让后地方力量建立台湾民主国并抵抗日军登陆。核心取证：《清德宗实录》光绪二十一年相关条；军机处档、战事方略及地方奏报。编年证词定位国家叙述中的时间，却需同地方或对方材料校验范围。来源保留：以光绪二十一年（1895年）的同期文书为定位起点；官修记录、奏报或外交文本服务特定机构立场，具体日期、规模、地方执行与对手视角须以独立材料复核。
\eventanchor{EQ-1895-TRIPLE-INTERVENTION}{三国干涉还辽改变战后列强竞争和清…}光绪二十一年（1895年），三国干涉还辽改变战后列强竞争和清廷外交处境。核心取证：《清德宗实录》光绪二十一年相关条；《筹办夷务始末》相关年卷与中外照会、条约文本。此类文件确认机构作出或收到的行政动作，不能跳过地方执行与反馈。来源保留：以光绪二十一年（1895年）的同期文书为定位起点；官修记录、奏报或外交文本服务特定机构立场，具体日期、规模、地方执行与对手视角须以独立材料复核。
\eventanchor{EQ-1896-LI-LOBANOV-TREATY}{《中俄密约》及东清铁路安排加强俄…}光绪二十二年（1896年），《中俄密约》及东清铁路安排加强俄国在东北的影响。核心取证：《清德宗实录》光绪二十二年相关条；《筹办夷务始末》相关年卷与中外照会、条约文本。此类文件确认机构作出或收到的行政动作，不能跳过地方执行与反馈。来源保留：以光绪二十二年（1896年）的同期文书为定位起点；官修记录、奏报或外交文本服务特定机构立场，具体日期、规模、地方执行与对手视角须以独立材料复核。
\eventanchor{EQ-1896-MODERN-EDUCATION-EXPANSION}{甲午战后新式学堂、译书与留学讨论…}光绪二十二年（1896年），甲午战后新式学堂、译书与留学讨论加速。核心取证：《清德宗实录》光绪二十二年相关条；内阁档、文集或报刊相关记录。编年证词定位国家叙述中的时间，却需同地方或对方材料校验范围。来源保留：以光绪二十二年（1896年）的同期文书为定位起点；官修记录、奏报或外交文本服务特定机构立场，具体日期、规模、地方执行与对手视角须以独立材料复核。
\eventanchor{EQ-1897-JIAOZHOU-BAY-OCCUPATION}{德国以巨野教案为由占领胶州湾，清…}光绪二十三年（1897年），德国以巨野教案为由占领胶州湾，清廷面临新的租借压力。核心取证：《清德宗实录》光绪二十三年相关条；《筹办夷务始末》相关年卷与中外照会、条约文本。此类文件确认机构作出或收到的行政动作，不能跳过地方执行与反馈。来源保留：以光绪二十三年（1897年）的同期文书为定位起点；官修记录、奏报或外交文本服务特定机构立场，具体日期、规模、地方执行与对手视角须以独立材料复核。
\eventanchor{EQ-1898-LEASED-TERRITORIES}{俄、英、法、德相继取得租借地或势…}光绪二十四年（1898年），俄、英、法、德相继取得租借地或势力范围安排。核心取证：《清德宗实录》光绪二十四年相关条；《筹办夷务始末》相关年卷与中外照会、条约文本。此类文件确认机构作出或收到的行政动作，不能跳过地方执行与反馈。来源保留：以光绪二十四年（1898年）的同期文书为定位起点；官修记录、奏报或外交文本服务特定机构立场，具体日期、规模、地方执行与对手视角须以独立材料复核。
\eventanchor{EQ-1898-HUNDRED-DAYS-REFORM}{光绪帝颁布变法诏令，推动官制、教…}光绪二十四年（1898年），光绪帝颁布变法诏令，推动官制、教育、经济与军事改革。核心取证：《清德宗实录》光绪二十四年相关条；宫中朱批奏折、内阁大库题本与《清史稿》相关志传。此类文件确认机构作出或收到的行政动作，不能跳过地方执行与反馈。来源保留：以光绪二十四年（1898年）的同期文书为定位起点；官修记录、奏报或外交文本服务特定机构立场，具体日期、规模、地方执行与对手视角须以独立材料复核。
\eventanchor{EQ-1898-WUXU-COUP}{慈禧发动政变，光绪帝被幽禁，戊戌…}光绪二十四年（1898年），慈禧发动政变，光绪帝被幽禁，戊戌变法主要措施被撤销或改写。核心取证：《清德宗实录》光绪二十四年相关条；宫中朱批奏折、内阁大库题本与《清史稿》相关志传。此类文件确认机构作出或收到的行政动作，不能跳过地方执行与反馈。来源保留：以光绪二十四年（1898年）的同期文书为定位起点；官修记录、奏报或外交文本服务特定机构立场，具体日期、规模、地方执行与对手视角须以独立材料复核。
\eventanchor{EQ-1899-BOXER-UPRISING}{义和团在华北扩展，反教、反洋冲突…}光绪二十五年（1899年），义和团在华北扩展，反教、反洋冲突加剧。核心取证：《清德宗实录》光绪二十五年相关条；地方志、督抚奏折及地方档案。编年证词定位国家叙述中的时间，却需同地方或对方材料校验范围。来源保留：以光绪二十五年（1899年）的同期文书为定位起点；官修记录、奏报或外交文本服务特定机构立场，具体日期、规模、地方执行与对手视角须以独立材料复核。
\subsection{1900—1909}
\eventanchor{EQ-1900-BOXER-WAR}{慈禧对列强宣战，八国联军入侵并占…}光绪二十六年（1900年），慈禧对列强宣战，八国联军入侵并占领北京。核心取证：《清德宗实录》光绪二十六年相关条；军机处档、战事方略及地方奏报。编年证词定位国家叙述中的时间，却需同地方或对方材料校验范围。来源保留：以光绪二十六年（1900年）的同期文书为定位起点；官修记录、奏报或外交文本服务特定机构立场，具体日期、规模、地方执行与对手视角须以独立材料复核。
\eventanchor{EQ-1900-SOUTHEAST-MUTUAL-PROTECTION}{东南互保使部分督抚与列强协商维持…}光绪二十六年（1900年），东南互保使部分督抚与列强协商维持地方秩序。核心取证：《清德宗实录》光绪二十六年相关条；宫中朱批奏折、内阁大库题本与《清史稿》相关志传。此类文件确认机构作出或收到的行政动作，不能跳过地方执行与反馈。来源保留：以光绪二十六年（1900年）的同期文书为定位起点；官修记录、奏报或外交文本服务特定机构立场，具体日期、规模、地方执行与对手视角须以独立材料复核。
\eventanchor{EQ-1901-BOXER-PROTOCOL}{《辛丑条约》签订，赔款、驻军和外…}光绪二十七年（1901年），《辛丑条约》签订，赔款、驻军和外交条件进一步限制清廷。核心取证：《清德宗实录》光绪二十七年相关条；《筹办夷务始末》相关年卷与中外照会、条约文本。此类文件确认机构作出或收到的行政动作，不能跳过地方执行与反馈。来源保留：以光绪二十七年（1901年）的同期文书为定位起点；官修记录、奏报或外交文本服务特定机构立场，具体日期、规模、地方执行与对手视角须以独立材料复核。
\eventanchor{EQ-1901-ZONGLI-YAMEN-REPLACED}{总理衙门改为外务部，晚清外交机构…}光绪二十七年（1901年），总理衙门改为外务部，晚清外交机构进一步制度化。核心取证：《清德宗实录》光绪二十七年相关条；宫中朱批奏折、内阁大库题本与《清史稿》相关志传。此类文件确认机构作出或收到的行政动作，不能跳过地方执行与反馈。来源保留：以光绪二十七年（1901年）的同期文书为定位起点；官修记录、奏报或外交文本服务特定机构立场，具体日期、规模、地方执行与对手视角须以独立材料复核。
\eventanchor{EQ-1902-CIXI-RETURN-BEIJING}{慈禧与光绪帝返回北京，战后朝廷开…}光绪二十八年（1902年），慈禧与光绪帝返回北京，战后朝廷开始推出新政措施。核心取证：《清德宗实录》光绪二十八年相关条；宫中朱批奏折、内阁大库题本与《清史稿》相关志传。此类文件确认机构作出或收到的行政动作，不能跳过地方执行与反馈。来源保留：以光绪二十八年（1902年）的同期文书为定位起点；官修记录、奏报或外交文本服务特定机构立场，具体日期、规模、地方执行与对手视角须以独立材料复核。
\eventanchor{EQ-1902-NEW-POLICIES-EDUCATION}{清末新政在学制、军制、商政和地方…}光绪二十八年（1902年），清末新政在学制、军制、商政和地方行政方面逐步启动。核心取证：《清德宗实录》光绪二十八年相关条；内阁档、文集或报刊相关记录。编年证词定位国家叙述中的时间，却需同地方或对方材料校验范围。来源保留：以光绪二十八年（1902年）的同期文书为定位起点；官修记录、奏报或外交文本服务特定机构立场，具体日期、规模、地方执行与对手视角须以独立材料复核。
\eventanchor{EQ-1903-BANNER-POSTS-REFORM}{清廷废除部分旗人官职垄断，调整旗…}光绪二十九年（1903年），清廷废除部分旗人官职垄断，调整旗籍制度与职务资格。核心取证：《清德宗实录》光绪二十九年相关条；地方志、督抚奏折及地方档案。编年证词定位国家叙述中的时间，却需同地方或对方材料校验范围。来源保留：以光绪二十九年（1903年）的同期文书为定位起点；官修记录、奏报或外交文本服务特定机构立场，具体日期、规模、地方执行与对手视角须以独立材料复核。
\eventanchor{EQ-1903-SHANGHAI-COMMERCIAL-LAW}{商法、公司与工商管理的近代制度讨…}光绪二十九年（1903年），商法、公司与工商管理的近代制度讨论在口岸和中央之间推进。核心取证：《清德宗实录》光绪二十九年相关条；户部奏销、盐政、河工或海关档案。此类文件确认机构作出或收到的行政动作，不能跳过地方执行与反馈。来源保留：以光绪二十九年（1903年）的同期文书为定位起点；官修记录、奏报或外交文本服务特定机构立场，具体日期、规模、地方执行与对手视角须以独立材料复核。
\eventanchor{EQ-1904-RUSSO-JAPANESE-WAR}{日俄战争在中国东北爆发，清廷名义…}光绪三十年（1904年），日俄战争在中国东北爆发，清廷名义中立而东北遭受战场化。核心取证：《清德宗实录》光绪三十年相关条；军机处档、战事方略及地方奏报。编年证词定位国家叙述中的时间，却需同地方或对方材料校验范围。来源保留：以光绪三十年（1904年）的同期文书为定位起点；官修记录、奏报或外交文本服务特定机构立场，具体日期、规模、地方执行与对手视角须以独立材料复核。
\eventanchor{EQ-1904-IMPERIAL-EXAMINATION-DEBATE}{科举存废与新式教育的争论加速}光绪三十年（1904年），科举存废与新式教育的争论加速。核心取证：《清德宗实录》光绪三十年相关条；内阁档、文集或报刊相关记录。编年证词定位国家叙述中的时间，却需同地方或对方材料校验范围。来源保留：以光绪三十年（1904年）的同期文书为定位起点；官修记录、奏报或外交文本服务特定机构立场，具体日期、规模、地方执行与对手视角须以独立材料复核。
\eventanchor{EQ-1905-ABOLITION-OF-CIVIL-SERVICE-EXAM}{清廷废除科举，延续千余年的取士制…}光绪三十一年（1905年），清廷废除科举，延续千余年的取士制度终结。核心取证：《清德宗实录》光绪三十一年相关条；内阁档、文集或报刊相关记录。编年证词定位国家叙述中的时间，却需同地方或对方材料校验范围。来源保留：以光绪三十一年（1905年）的同期文书为定位起点；官修记录、奏报或外交文本服务特定机构立场，具体日期、规模、地方执行与对手视角须以独立材料复核。
\eventanchor{EQ-1905-TONGMENGHUI}{中国同盟会在东京成立，革命组织和…}光绪三十一年（1905年），中国同盟会在东京成立，革命组织和宣传网络进一步整合。核心取证：《清德宗实录》光绪三十一年相关条；地方志、督抚奏折及地方档案。编年证词定位国家叙述中的时间，却需同地方或对方材料校验范围。来源保留：以光绪三十一年（1905年）的同期文书为定位起点；官修记录、奏报或外交文本服务特定机构立场，具体日期、规模、地方执行与对手视角须以独立材料复核。
\eventanchor{EQ-1906-CONSTITUTIONAL-PROMISE}{清廷宣布预备立宪，着手考察各国制…}光绪三十二年（1906年），清廷宣布预备立宪，着手考察各国制度和改革官制。核心取证：《清德宗实录》光绪三十二年相关条；宫中朱批奏折、内阁大库题本与《清史稿》相关志传。此类文件确认机构作出或收到的行政动作，不能跳过地方执行与反馈。来源保留：以光绪三十二年（1906年）的同期文书为定位起点；官修记录、奏报或外交文本服务特定机构立场，具体日期、规模、地方执行与对手视角须以独立材料复核。
\eventanchor{EQ-1906-NEW-ARMY-REFORM}{新军编练与军事学堂扩展，军队成为…}光绪三十二年（1906年），新军编练与军事学堂扩展，军队成为新式政治动员的重要空间。核心取证：《清德宗实录》光绪三十二年相关条；军机处档、战事方略及地方奏报。编年证词定位国家叙述中的时间，却需同地方或对方材料校验范围。来源保留：以光绪三十二年（1906年）的同期文书为定位起点；官修记录、奏报或外交文本服务特定机构立场，具体日期、规模、地方执行与对手视角须以独立材料复核。
\eventanchor{EQ-1907-MANCHURIAN-PROVINCES}{清廷将东三省调整为行省体制，加强…}光绪三十三年（1907年），清廷将东三省调整为行省体制，加强近代行政与边防经营。核心取证：《清德宗实录》光绪三十三年相关条；军机处档、理藩院档或相关方略。编年证词定位国家叙述中的时间，却需同地方或对方材料校验范围。来源保留：以光绪三十三年（1907年）的同期文书为定位起点；官修记录、奏报或外交文本服务特定机构立场，具体日期、规模、地方执行与对手视角须以独立材料复核。
\eventanchor{EQ-1907-QING-LEGAL-REFORM}{清廷推进修律、审判和警政改革，尝…}光绪三十三年（1907年），清廷推进修律、审判和警政改革，尝试改造传统法律制度。核心取证：《清德宗实录》光绪三十三年相关条；宫中朱批奏折、内阁大库题本与《清史稿》相关志传。此类文件确认机构作出或收到的行政动作，不能跳过地方执行与反馈。来源保留：以光绪三十三年（1907年）的同期文书为定位起点；官修记录、奏报或外交文本服务特定机构立场，具体日期、规模、地方执行与对手视角须以独立材料复核。
\eventanchor{EQ-1908-GUANGXU-CIXI-DEATHS}{光绪帝与慈禧太后相继去世，溥仪继…}光绪三十四年（1908年），光绪帝与慈禧太后相继去世，溥仪继位为宣统帝。核心取证：《清德宗实录》光绪三十四年相关条；宫中朱批奏折、内阁大库题本与《清史稿》相关志传。此类文件确认机构作出或收到的行政动作，不能跳过地方执行与反馈。来源保留：以光绪三十四年（1908年）的同期文书为定位起点；官修记录、奏报或外交文本服务特定机构立场，具体日期、规模、地方执行与对手视角须以独立材料复核。
\eventanchor{EQ-1908-CONSTITUTIONAL-OUTLINE}{清廷公布《钦定宪法大纲》等文件，…}光绪三十四年（1908年），清廷公布《钦定宪法大纲》等文件，明确预备立宪框架。核心取证：《清德宗实录》光绪三十四年相关条；宫中朱批奏折、内阁大库题本与《清史稿》相关志传。此类文件确认机构作出或收到的行政动作，不能跳过地方执行与反馈。来源保留：以光绪三十四年（1908年）的同期文书为定位起点；官修记录、奏报或外交文本服务特定机构立场，具体日期、规模、地方执行与对手视角须以独立材料复核。
